\section*{Supplementary material}
\addcontentsline{toc}{section}{Supplementary material}

\subsection*{S1. Full BH-FDR hypothesis family (17 tests)}

Table~\ref{tab:supp-bhfdr} lists, in full, the Benjamini--Hochberg false-discovery-rate (BH-FDR)
family described in Section~\ref{sec:gates}: 13 marker hypotheses against real ground truth
plus four nested/refit incremental-value tests from the confounder audit. The marker rows are
each a genuinely different clinical or molecular
target tested with CONCH embeddings under one fixed protocol (patient-disjoint 5-fold grouped
cross-validation; standardized-and-ridge-regularized regression for continuous targets;
standardized, balanced-class logistic regression for binary targets -- full specification in
Methods). The six initial frozen-pool markers (rows 1--6) appear in the main text
(Table~\ref{tab:pool}) with their patient-level effect size and $q$-value only; here we also
report the uncorrected patient-level $p$-value alongside $q$, and include the seven additional
TCGA-PRAD molecular candidates (rows 7--13) that were tested but not promoted to the marker
pool. ETV4 is no longer below 0.05 after expanding the frozen family ($q=0.056$) and was in any
case pre-designated as excluded from both pools. Rows 14--17 are the nested audits.

\textbf{Family membership.} All 17 rows below share one BH-FDR family (Benjamini--Hochberg
applied jointly). \textbf{Excluded} from this family, as
already noted in Section~\ref{sec:gates}: external-cohort re-tests of an already-tested
hypothesis (PANDA zero-shot transfer of markers 1--2, PRECISE real-Gleason validation of marker
1, marker 7's TCGA-PRAD and PFS/DFS transfer tests) and cross-encoder (Virchow) replications of
any of the 13 rows below. These are confirmatory re-evaluations of a hypothesis already counted
once in the family, not new comparisons that inflate the family-wise false-discovery risk the
way testing 13 distinct candidate markers does; correcting across both the original tests and
their confirmatory replications would double-count the same hypothesis and over-penalize the
family.

\begingroup
\footnotesize
\begin{longtable}{@{}clrp{3.6cm}rr@{}}
\toprule
\# & Test & $n$ & Patient-level effect [95\% CI] & Uncorr.\ $p$ & $q$ \\
\midrule
\endhead
1 & NADT H\&E$\to$Gleason ($\rho$) & 39 & $+0.478$ $[+0.170,+0.712]$ & $2.1\times10^{-3}$ & $7.0\times10^{-3}$ \\
2 & NADT H\&E$\to$Phenotype ($\rho$) & 39 & $+0.800$ $[+0.625,+0.909]$ & $9.6\times10^{-10}$ & $1.6\times10^{-8}$ \\
3 & NADT ERG$\to$Gleason ($\rho$) & 39 & $+0.524$ $[+0.216,+0.749]$ & $6.1\times10^{-4}$ & $3.5\times10^{-3}$ \\
4 & TCGA-PRAD H\&E$\to$PTEN loss (AUROC) & 273 & $0.632$ $[0.560,0.706]$ & $5.6\times10^{-4}$ & $3.5\times10^{-3}$ \\
5 & TCGA-PRAD H\&E$\to$SPOP mut.\ (AUROC) & 273 & $0.519$ $[0.408,0.627]$ & $0.742$ & $0.788$ \\
6 & TCGA-PRAD H\&E$\to$AR score ($\rho$) & 273 & $+0.195$ $[+0.074,+0.310]$ & $1.2\times10^{-3}$ & $5.2\times10^{-3}$ \\
7 & TCGA-PRAD H\&E$\to$TP53 mut.\ (AUROC) & 273 & $0.485$ $[0.354,0.614]$ & $0.835$ & $0.835$ \\
8 & TCGA-PRAD H\&E$\to$TP53 CNA loss (AUROC) & 273 & $0.562$ $[0.492,0.632]$ & $0.098$ & $0.208$ \\
9 & TCGA-PRAD H\&E$\to$RB1 CNA loss (AUROC) & 273 & $0.536$ $[0.453,0.615]$ & $0.363$ & $0.441$ \\
10 & TCGA-PRAD H\&E$\to$SPINK1 high (AUROC) & 273 & $0.618$ $[0.470,0.762]$ & $0.152$ & $0.288$ \\
11 & TCGA-PRAD H\&E$\to$ETV1 altered (AUROC) & 273 & $0.429$ $[0.307,0.547]$ & $0.245$ & $0.347$ \\
12 & TCGA-PRAD H\&E$\to$ETV4 altered (AUROC) & 273 & $0.687$ $[0.517,0.843]$ & $0.023$ & $0.056$ \\
13 & TCGA-PRAD H\&E$\to$ERG fusion (AUROC) & 273 & $0.526$ $[0.456,0.592]$ & $0.458$ & $0.519$ \\
14 & Nested PTEN increment ($\Delta$AUROC) & 273 & $+0.019$ $[-0.025,+0.068]$ & $0.203$ & $0.314$ \\
15 & Nested AR increment ($\Delta R^2$) & 273 & $+0.004$ $[-0.036,+0.042]$ & $0.176$ & $0.299$ \\
16 & Nested marker 7, grade-only ($\Delta C$) & 270 & $+0.062$ $[+0.008,+0.119]$ & $0.0100$ & $0.028$ \\
17 & Nested marker 7, fully adjusted ($\Delta C$) & 153 & $+0.002$ $[-0.012,+0.015]$ & $0.305$ & $0.399$ \\
\bottomrule
\caption{All 17 tests in the BH-FDR family, patient-level, with uncorrected $p$ and
BH-FDR $q$ shown side by side. Rows 1--6 are the initially frozen pool markers (numbered as in
the main text); rows 7--13 are additional TCGA-PRAD molecular candidates tested but not
promoted to the pool; rows 14--17 are held-out increments. Effect size is Spearman $\rho$ for
continuous targets, AUROC for binary targets, and the named paired increment for audit rows;
95\% CIs are patient-ID bootstrap (2{,}000 resamples).}
\label{tab:supp-bhfdr}
\end{longtable}
\endgroup

\subsection*{S2. Fully-adjusted clinical covariate model for marker 7 (recurrence)}

Table~\ref{tab:supp-marker7-clinical} extends the marker-7 confounder audit
(Section~\ref{sec:res-marker7-clinical}, Table~\ref{tab:confounder-7}) beyond Gleason grade
alone, addressing a reviewer question about whether additional TCGA clinical covariates (PSA,
T-stage, age, surgical margin status) change the conclusion. See Section~\ref{sec:res-marker7-clinical}
for the discussion of this result.

\begin{table}[htbp]
\centering
\small
\begin{tabular}{p{7.2cm}ccc}
\toprule
Model & $n$ & Held-out C-index & Refit permutation $p$ \\
\midrule
Clinical-only (grade only) & 270 & 0.645 & -- \\
Combined (grade + marker7) & 270 & 0.707 & 0.0100 \\
Fully-adjusted clinical-only (grade + age + T-stage + PSA + margin) & 153 & 0.839 & -- \\
Fully-adjusted combined (fully-adjusted clinical + marker7) & 153 & 0.840 & 0.305 \\
M4 clinical + site (regularized sensitivity) & 153 & 0.881 & -- \\
M5 M4 + marker7 (regularized sensitivity) & 153 & 0.879 & -- \\
\bottomrule
\end{tabular}
\caption{Nested marker 7 audit, grade-only vs.\ fully-adjusted clinical baseline. Rows 1--2 and
3--4 correspond to the grade-only and fully-adjusted $\Delta$ values in
Table~\ref{tab:confounder-7}, shown here as clinical/combined pairs rather than as increments.
The fully-adjusted rows use a
complete-case subset ($n=153$ of the 270-patient marker-7 cohort; PSA is the limiting
covariate, populated for only 167/270 patients in cBioPortal \texttt{prad\_tcga\_pub}) and add
age, ordinal T-stage (major stage 1--4 only), PSA, and binary surgical margin status
(positive = R1/R2). Surgical margin status alone is a very strong predictor in this cohort
(recurrence-event rate 67\% among margin-positive vs.\ 7\% among margin-negative patients,
hazard ratio $\approx 13$). These historical conditional point estimates are superseded for
comparative interpretation by the completed same-patient/same-draw R4 analysis on the same
153-patient complete-case cohort; it does not establish robust independent or incremental
prognostic value.}
\label{tab:supp-marker7-clinical}
\end{table}

\subsection*{S3. Completed Gate A stability grid and quality control}

The completed frozen audit contains 360 correlated seed-level cells spanning six markers,
CONCH and Virchow, 16/32/64 tiles per slide, five tile-sampling seeds, and native versus shared
$1.76\,\mu$m/px scales. The saved 72-row configuration summary and 390-row paired-contrast
table are the sole numeric sources for Figure~\ref{fig:scale-heatmap} and the table below.
Configuration Student-$t$ intervals use five sampling seeds (df=4) on the same frozen cohorts
and folds; they are explicitly not patient or population confidence intervals. Contrast rows
preserve pair identity and report descriptive null crossings without independence-based
$p$-values.

\begin{table}[htbp]
\centering
\scriptsize
\resizebox{\textwidth}{!}{%
\begin{tabular}{llccccc}
\toprule
Marker & Metric & Observed cell range & Chance-or-worse & Seed-null straddle & Native-scale crossing & Shared-scale encoder crossing \\
\midrule
Gleason & Spearman $\rho$ & 0.271--0.646 & 0/60 & 0/12 & 0/30 & 0/15 \\
Phenotype & AUROC & 0.845--0.965 & 0/60 & 0/12 & 0/30 & 0/15 \\
PTEN & AUROC & 0.519--0.717 & 0/60 & 0/12 & 0/30 & 0/15 \\
SPOP & AUROC & 0.348--0.679 & 25/60 & 8/12 & 19/30 & 4/15 \\
AR & Spearman $\rho$ & 0.136--0.297 & 0/60 & 0/12 & 0/30 & 0/15 \\
Marker 7 & C-index & 0.481--0.755 & 1/60 & 1/12 & 1/30 & 0/15 \\
\bottomrule
\end{tabular}%
}
\caption{Frozen Gate A stability-grid summary derived from the saved 72-configuration and 390-contrast CSVs. Observed ranges span the 60 seed-level cells per marker. Seed-null straddles summarize 12 five-seed configurations; sampling seeds reuse the same frozen cohorts and folds and are not patient-level confidence intervals.}
\label{tab:supp-stability-grid}
\end{table}


The aggregation reconciled 360 cell rows, 1{,}800 fold rows, and 60 coordinate shards before
publishing derived files. Overall, 26/360 cells were chance-or-worse, 20/180 native-scale pairs
crossed the marker-specific null, 4/90 shared-scale encoder pairs crossed it, and 9/72
configuration seed ranges straddled the null. Marker-7 source retention varied from 498 to 502
patients; the common intersection was 498. R2 refit and rescored all 60 marker-7 cells with
that fixed 498-patient/85-event source and the unchanged 270-patient/57-event target. The
saved 60-row cell table, 65-row paired-delta table, and 10{,}040-row membership manifest record
zero raw-to-common null crossings, preserved directions for all 30 scale, 20 tile-budget, and
15 shared-scale encoder contrasts, and common-minus-raw shifts from $-0.0023$ to $+0.0229$.
This source-constant result is still a correlated sensitivity analysis, not a causal scale
effect, universal encoder comparison, or independent replication. Historical append/resume
logs lack complete per-run command/config manifests; that limitation is retained rather than
reconstructed retrospectively.

The NADT input inventory underwent a complete 463/463 TIFF header scan. Two slides emitted an
invalid-page-offset header warning; both remained represented in all expected coordinate
shards. This is a header-level QC observation only and does not establish whether pixels were
or were not affected.

R5 saves all 22 AR site descriptors, all 300 page-0 metadata rows, all 22 SPOP site rows, two
power rows, and 433 reconstructed SPOP leave-one-site-out predictions. The AR forest-eligible
frame is 224 slides/209 patients across CH, EJ, G9, HC, KK, and YL. Its site descriptors are
patient-level, but forest correlations are slide-level with patient-cluster bootstrap
intervals; all six intervals include zero. ScanScope raw IDs are present for 280/300 slides,
whereas explicit stain fields are present for 0/300. SPOP CH is single-class undefined; the
other five site intervals include or touch 0.5, and undefined bootstrap draws are retained.
The 80\%/90\% MDE values 0.661/0.685 are fixed-score normal approximations and not
effect-exclusion bounds.

\textbf{Historical scale-control bug context.} Before the frozen grid, an obsolete diagnostic
excluded the finest tiled pyramid level, causing distinct target scales to resolve to the same
page. The corrected frozen runners retain the finest level as a candidate and use windowed
physical-size reads. The abandoned 90-slide, three-marker, single-seed output is not used in
the current figure, table, or claim matrix.

\subsection*{S4. TCGA-PRAD recurrence-label provenance and the strict-endpoint sensitivity check}

The GDC-\texttt{follow\_ups}-derived recurrence label used throughout this paper (event =
\texttt{disease\_response} recorded as ``wt-with tumor'' at any follow-up visit) is generated by
a script that also exports a 500-case, patient-level provenance table: every raw follow-up
record considered, the specific field value that drove the event/censoring decision, and the
final label. The label-generating script and its output were re-verified against a
previously-cached copy by hash comparison and found identical (493 valid labels, 111 events),
confirming the label itself did not silently change between the original construction and this
paper's later analyses (Section~\ref{sec:res-label-benchmark}).

The strict, persistent-disease-excluding sensitivity endpoint discussed in
Section~\ref{sec:res-label-benchmark} applies one additional rule to this same provenance
table: an event only counts if the patient has a \emph{documented tumor-free} follow-up visit
strictly before a later with-tumor observation. Of 493 labeled cases (111 events under the
original, more permissive definition), only 11 satisfy this stricter rule; the other 100
with-tumor observations are a patient's first or only recorded assessment (or occur on the same
day as an earlier one), which this stricter rule treats as persistent/indeterminate disease
rather than as recurrence, and excludes rather than recodes as censored. Within the 270-patient
embedded cohort this leaves $n=220$ and only 7 events -- far too few for the C-index and
td-AUROC reported in Section~\ref{sec:res-label-benchmark} to be read as more than directionally
suggestive; we report the numbers because the direction is consistent with the primary result,
not because the sample size supports a standalone claim.

\subsection*{S5. Other scope-bounding decisions}

Two further deliberate scope reductions, noted here rather than in the main text since they
affect \emph{how much} evidence a check provides rather than \emph{what} it found: (1) the
Virchow cross-check of the nested confounder audit (Section~\ref{sec:confounder-audits}) was
run only for PTEN and AR, not marker 7. Gate A now treats marker 7 as encoder-by-scale
conditioned. Its common-498 R2, official-PFI R3, and paired common-cohort R4 analyses are
complete and retain an endpoint-specific exploratory interpretation. (2) The M0--M5 clinical
hierarchy (Section~\ref{sec:res-marker7-clinical}) uses complete-case handling throughout;
multiple imputation was not performed because correctly combining it with nested cross-fitting
would require a separate approved analysis.
